% =====================================================================
%  Section 4: Experiments
%  Main setting: LLaMA-2-7B full fine-tuning, Alpaca-GPT4,
%  10% selection, nine evaluation tasks across five capability groups.
% =====================================================================
\section{Experiments}
\label{sec:experiments}

We evaluate \textsc{TADS} as a dynamic re-ranking method for
instruction-tuning data selection. The main setting uses
LLaMA-2-7B full fine-tuning on Alpaca-GPT4 with a $10\%$ selection
budget and a nine-task evaluation suite spanning knowledge,
reasoning, code, and multilingual understanding. This setting
allows a controlled comparison against static data-selection
methods, a no-anchor composite-reward base, and full-data
fine-tuning under a shared training and evaluation budget.

The experiments answer four questions. \textbf{RQ1}: Does
\textsc{TADS} improve over static and dynamic data-selection
baselines? \textbf{RQ2}: How sensitive is performance to the
alignment weight $\lambda$? \textbf{RQ3}: Does the trajectory
anchor satisfy the stability bound in real training? \textbf{RQ4}:
How much re-ranking overhead does the structural signal add? A
larger Qwen2.5-14B + LoRA scaling study and broader extensions are
reported in App.~\ref{app:scaling} and App.~\ref{app:inprogress}.

\subsection{Experimental setup}
\label{sec:setup}

\paragraph{Model and fine-tuning.}
The main experiments use LLaMA-2-7B~\cite{touvron2023llama2} with
\emph{full fine-tuning}: all parameters are updated, with no
low-rank adapter. All methods share the same optimiser, learning
rate schedule, batch size, number of epochs, and selection budget,
so that the selected data is the only varying factor.

\paragraph{Candidate pool and selection budget.}
The candidate pool is Alpaca-GPT4~\cite{alpaca_gpt4}, a
self-instruct instruction-tuning corpus synthesised by GPT-4. The
selection budget is fixed at $\rho=10\%$. The full-data reference
uses the entire pool and serves as the $100\%$ fine-tuning baseline.

\paragraph{Evaluation.}
We evaluate on nine tasks grouped into five capability categories:
\emph{factual knowledge} (MMLU, MMLU-Pro), \emph{mathematical
reasoning} (GSM8K, SVAMP), \emph{coding ability} (HumanEval, MBPP),
\emph{multilingual understanding} (TydiQA, XQuAD), and
\emph{general reasoning} (BBH). Per-task scoring conventions and
worked examples are listed in App.~\ref{app:examples}. We report
per-task scores, the nine-task average (AVG), and the relative
change $\Delta$ over the Alpaca-GPT4 full-data baseline.

\paragraph{Baselines.}
We compare \textsc{TADS} against three groups of baselines.
\emph{(i) Data-side references}: Alpaca-GPT4 full fine-tuning
($100\%$) and a Random $10\%$ subset. \emph{(ii) Static IT-data
selection methods}: LIMA~\cite{zhou2023lima},
AlpaGasus~\cite{alpagasus}, Q2Q~\cite{q2q},
SelectIT~\cite{selectit}, LESS~\cite{less}, and
NAIT~\cite{nait2026}. LESS scores are cited because a full re-run
exceeds our compute budget; its reported results cover only a
five-task subset, so the remaining cells are blank in
Table~\ref{tab:main}. \emph{(iii) Dynamic re-ranking}: a
composite-reward dynamic selector~\cite{kung2023activeit,
sun2025dots} ranking candidates by the composite reward
$\mathcal{R}_i^{(t)} = w^{(t)}\mathcal{L}_i^{(t)} +
(1-w^{(t)})\mathcal{H}_i^{(t)}$ (Eq.~\ref{eq:composite}). The
$\lambda=0$ setting removes the trajectory anchor and recovers
this composite-reward base ranking; the comparison with
\textsc{TADS}($\lambda{=}1$) therefore isolates the effect of the
anchor.

\subsection{Main result}
\label{sec:main}

Table~\ref{tab:main} reports the main comparison. The final two
rows give the controlled dynamic re-ranking comparison: the base
uses composite reward $\mathcal{R}_i^{(t)}$ with
$\lambda=0$, while \textsc{TADS} uses the same pipeline with the
trajectory anchor enabled at $\lambda=1$. Any difference between
these rows is attributable to the trajectory anchor, not to
optimisation, selection ratio, or evaluation protocol.

\begin{table*}[t]
\centering
\caption{\textbf{Main result on LLaMA-2-7B, full fine-tuning,
Alpaca-GPT4, $10\%$ selection.} Nine tasks are grouped into five
capability families; AVG is the nine-task mean and $\Delta$ is the
relative change over the Alpaca-GPT4 full-data baseline. The
$\lambda=0$ row is the composite-reward base ranking
$\mathcal{R}_i^{(t)}$ (Eq.~\ref{eq:composite}); the
\textsc{TADS}-vs-base comparison isolates the trajectory anchor.
Best per column in \textbf{bold}. \STATUS{Dynamic-selection rows
are preliminary; MMLU-Pro, H-Eval, and AVG/$\Delta$ pending the
remaining evaluations. Non-dynamic rows are carried-over reference
values and are shown with strikethrough while the full comparison
is being finalised.}}
\label{tab:main}
\small
\setlength{\tabcolsep}{4pt}
\begin{tabular}{llcccccccccc}
\toprule
 & & \multicolumn{2}{c}{Factual Know.}
   & \multicolumn{2}{c}{Math. Reason.}
   & \multicolumn{2}{c}{Coding}
   & \multicolumn{2}{c}{Multilingual}
   & Gen. R. & \\
\cmidrule(lr){3-4}\cmidrule(lr){5-6}\cmidrule(lr){7-8}\cmidrule(lr){9-10}\cmidrule(lr){11-11}
ID & Method & MMLU & MMLU-Pro & GSM & SVAMP & H-Eval & MBPP
   & TydiQA & XQuAD & BBH & AVG / $\Delta$ \\
\midrule
\multicolumn{12}{l}{\textit{Data-side references}}\\
01 & Alpaca-GPT4 (Full)
& \sout{46.87} & \sout{21.89} & \sout{14.63} & \sout{39.00}
& \sout{27.87} & \sout{51.58} & \sout{39.48} & \sout{42.99}
& \sout{39.94} & \sout{36.03} / --- \\

06 & Random ($10\%$)
& \sout{47.14} & \sout{21.43} & \sout{14.13} & \sout{35.67}
& \sout{25.55} & \sout{47.35} & \sout{44.16} & \sout{46.56}
& \sout{39.21} & \sout{35.69} / \sout{$-0.94\%$} \\

\midrule
\multicolumn{12}{l}{\textit{Static IT-data selection}}\\
02 & LIMA
& \sout{45.20} & \sout{23.04} & \sout{15.76} & \sout{37.67}
& \sout{27.75} & \sout{46.56} & \sout{44.92} & \sout{44.72}
& \sout{39.91} & \sout{36.17} / \sout{$+0.39\%$} \\

03 & AlpaGasus
& \sout{43.21} & \sout{21.96} & \sout{13.34} & \sout{36.67}
& \sout{23.94} & \sout{46.08} & \sout{44.70} & \sout{46.84}
& \sout{39.91} & \sout{35.18} / \sout{$-2.34\%$} \\

04 & Q2Q
& \sout{46.73} & \sout{21.50} & \sout{14.50} & \sout{35.00}
& \sout{25.19} & \sout{44.97} & \sout{44.41} & \sout{48.44}
& \sout{40.34} & \sout{35.68} / \sout{$-0.98\%$} \\

05 & SelectIT
& \sout{47.90} & \sout{22.86} & \sout{15.40} & \sout{41.11}
& \sout{27.92} & \sout{49.47} & \sout{43.91} & \sout{45.56}
& \sout{40.33} & \sout{37.16} / \sout{$+3.15\%$} \\

10 & LESS$^{\dagger}$
& \sout{48.12} & \sout{24.50} & -- & --
& -- & -- & \sout{43.42} & \sout{45.51}
& \sout{38.15} & -- / -- \\

07 & NAIT (All)
& \sout{46.83} & \sout{23.29} & \sout{16.53} & \sout{39.67}
& \sout{26.44} & \sout{47.62} & \sout{46.09} & \sout{48.27}
& \sout{40.02} & \sout{37.20} / \sout{$+3.24\%$} \\

\midrule
\midrule
\multicolumn{12}{l}{\textit{Dynamic re-ranking}}\\
08 & Composite-reward base ($\lambda=0$)
& 45.90 & \STATUS{??.??} & 16.48 & 37.67
& \STATUS{??.??} & 46.58 & 40.61 & 46.76
& 40.21 & \STATUS{??.??} / \STATUS{?} \\

09 & \textbf{TADS} ($\lambda=1$)
& 47.17 & \STATUS{??.??} & 16.63 & 39.03
& \STATUS{??.??} & 47.35 & 41.64 & 46.33
& 40.36 & \STATUS{??.??} / \STATUS{?} \\

\bottomrule
\end{tabular}
\end{table*}

\paragraph{Effect of the trajectory anchor.}
Rows 08--09 form the cleanest test of the proposed mechanism: both
use the same composite-reward base ranking
$\mathcal{R}_i^{(t)}$, and \textsc{TADS} differs only by
the multiplicative trajectory-anchor factor. The completed metrics
already show consistent improvements for \textsc{TADS} on MMLU,
GSM8K, SVAMP, MBPP, TydiQA, and BBH, while the remaining
evaluations are still pending. This supports the central claim
that a trajectory-derived structural signal can improve a dynamic
re-ranker without requiring seed examples, validation labels, or
auxiliary scorers.

\paragraph{Comparison to static selectors.}
The strongest static baselines are SelectIT and NAIT (All), both
of which exceed the full-data reference on the nine-task average.
\textsc{TADS} differs from these methods in the source and timing
of the selection signal: rather than using an external judge, a
target set, or a seed-derived capability direction fixed before
training, it refreshes a model-intrinsic trajectory anchor during
training. LESS is strong on the factual-knowledge tasks it
targets, but its reported subset does not support a full nine-task
average.

\paragraph{Selection can beat scale.}
The table also reflects a recurring observation in instruction
tuning: several $10\%$ selection methods match or outperform
full-data fine-tuning. This motivates selection criteria that are
not only accurate but also adaptive, stable, and inexpensive.

\subsection{$\lambda$-ablation of the alignment weight}
\label{sec:lambda-ablation}

We validate the core \textsc{TADS} lever, the alignment weight
$\lambda$ in Eq.~\ref{eq:score}. Table~\ref{tab:lambda-7b} reports
the main LLaMA-2-7B sweep, while Table~\ref{tab:lambda-05b}
reports the Qwen2.5-0.5B sanity-track sweep. In both settings,
$\lambda=0$ recovers the composite-reward base ranking, and
larger values strengthen the trajectory-anchor factor.

The expected pattern is non-monotonic. Small non-zero $\lambda$
adds structural signal, whereas overly large values saturate the
multiplicative boost
$1+\lambda\,\widetilde{\mathrm{align}}_i^{(t)} \in [1,1+\lambda]$.
The 0.5B sanity track already exhibits this concave behaviour,
peaking at $\lambda=0.1$.

\begin{table}[t]
\centering
\caption{$\lambda$-ablation of the alignment weight in
Eq.~\ref{eq:score}, on LLaMA-2-7B with full fine-tuning, Alpaca-GPT4
at $\rho=10\%$. $\lambda=0$ recovers the composite-reward base
ranking; larger $\lambda$ strengthens the trajectory-anchor signal.
Avg is the four-task mean. Best per column in \textbf{bold}.}
\label{tab:lambda-7b}
\begin{tabular}{lccccc}
\toprule
$\lambda$ & MMLU & GSM8K & H-Eval & TydiQA & Avg \\
\midrule
$0$ (Composite-reward base) & \STATUS{??.??} & \STATUS{??.??} & \STATUS{??.??} & \STATUS{??.??} & \STATUS{??.??} \\
$0.1$                & \STATUS{??.??} & \STATUS{??.??} & \STATUS{??.??} & \STATUS{??.??} & \STATUS{??.??} \\
$0.5$                & \STATUS{??.??} & \STATUS{??.??} & \STATUS{??.??} & \STATUS{??.??} & \STATUS{??.??} \\
$1.0$ (Default)      & \STATUS{??.??} & \STATUS{??.??} & \STATUS{??.??} & \STATUS{??.??} & \STATUS{??.??} \\
$5.0$                & \STATUS{??.??} & \STATUS{??.??} & \STATUS{??.??} & \STATUS{??.??} & \STATUS{??.??} \\
$10.0$               & \STATUS{??.??} & \STATUS{??.??} & \STATUS{??.??} & \STATUS{??.??} & \STATUS{??.??} \\
\bottomrule
\end{tabular}
\end{table}

\begin{table}[t]
\centering
\caption{$\lambda$-ablation on the Qwen2.5-0.5B sanity track,
fine-tuned on Alpaca-GPT4 at $\rho=10\%$ for 3 epochs.
$\lambda=0$ recovers the composite-reward base ranking;
larger $\lambda$ strengthens the trajectory-anchor signal. Avg
averages MMLU and GSM8K. Best per column in \textbf{bold}.}
\label{tab:lambda-05b}
\begin{tabular}{lccc}
\toprule
$\lambda$ & MMLU & GSM8K & Avg \\
\midrule
$0$ (Composite-reward base) & 46.09 & 28.28 & 37.18 \\
$0.1$                & 46.96 & \textbf{30.17} & \textbf{38.56} \\
$1.0$ (default)      & \textbf{47.03} & 28.96 & 37.99 \\
$5.0$                & 46.85 & 29.49 & 38.17 \\
$10.0$               & 46.90 & 28.73 & 37.81 \\
\bottomrule
\end{tabular}
\end{table}

% Red-border status box. Requires \usepackage{xcolor}.
\begin{center}
\fcolorbox{red}{white}{%
\begin{minipage}{0.95\columnwidth}
\color{red}\small
\textbf{TODO:}
\begin{itemize}\setlength\itemsep{0.05em}
  \item \textbf{Scale matrix}: 0.5B / 7B / 14B $\times$
    \{random\_10, consensus\_10, tads\_10, full\_100\} $\times$
    5 benchmarks. 0.5B sanity $\checkmark$, 7B running, 14B LoRA scheduled.
  \item \textbf{Dataset axis}: WizardLM Evol-Instruct 70K, Qwen2.5-7B,
    same 4 methods. random\_10 / full\_100 sealed; tads\_10 /
    consensus\_10 re-training (DDP broadcast fix).
  \item \textbf{$\rho$-sweep}: $\rho \in \{0.05, 0.10, 0.20, 0.50\}$
    at $\lambda{=}1$ on Qwen2.5-7B. $\rho{=}0.10$ from main matrix,
    $\rho{=}1.0$ from full\_100, three new cells.
  \item \textbf{$\lambda$-ablation at 7B}: extend
    Table~\ref{tab:lambda-05b} to Qwen2.5-7B; new
    $\lambda \in \{0.5, 2, 5, 10\}$ ($\lambda{=}0, 1$ reuse existing).
\end{itemize}
\end{minipage}}
\end{center}

\subsection{Empirical verification of Theorem~\ref{thm:stability}}
\label{sec:thm-empirical}

Theorem~\ref{thm:stability} predicts that the trajectory anchor
moves by at most $(2C_\Sigma/\gamma_{\min})\eta_t$ between refresh
points, provided the covariance drift is controlled, the top
eigenvalue gap is positive, and the PCA sign is consistently
calibrated. We verify these conditions during Qwen2.5-0.5B + LoRA
training on Alpaca-GPT4 by running a verifier every $25$ optimiser
steps on a fixed 2{,}000-sample probe. For each layer and refresh
point, the verifier records the eigengap, covariance drift, anchor
displacement, sign inner-product, and learning rate.

The assumptions hold in the robust form used by the proof. Across
$68$ post-warmup refresh points and $24$ decoder layers, the
eigengap is strictly positive on the monitored layers, with
$\gamma_{\min}=0.050$ after thresholding. The trimmed covariance
drift gives $\hat{C}_\Sigma=5.78\times 10^3$, and the measured
sign-flip rate is $3.6\%$, below the $5\%$ tolerance. Full
per-layer diagnostics appear in App.~\ref{app:thm-verification}.

At the constant operating point $\eta=2\times 10^{-4}$, the
per-step bound is $(2\hat{C}_\Sigma/\gamma_{\min}) \eta = 46.24$,
and the measured per-step displacement stays strictly below this
bound at all $1{,}496$ layer-refresh measurements with zero
violations, median tightness $1.46\times 10^{-3}$, worst
$7.4\times 10^{-3}$. The stronger version of the guarantee is the
cumulative tail bound
$\|v_l^{(t+m)} - v_l^{(t)}\| \le (2\hat{C}_\Sigma/\gamma_{\min})
\sum_{s=t}^{t+m-1} \eta_s$, which under a decaying schedule shrinks
to zero as $t \to T$. We test this directly on App.~F cell~B
(probe $2{,}000$, cosine schedule, $78$ refresh points): both the
measured cumulative drift to the final anchor
$\|v_l^{(t)} - v_l^{(T)}\|$ and the theoretical tail bound decay
monotonically toward zero, with the measured curve tracking the
bound's shape rather than saturating (Fig.~\ref{fig:thm-main}).
This is the $\varepsilon$-stationarity statement we asked the
theorem for: the trajectory anchor is not just per-step bounded, it
\emph{converges}.

\paragraph{Bounded vs unbounded: same diagnostic, two regimes.}
As an additional post-hoc diagnostic, we compute the same
trajectory-anchor displacement statistics on a matched no-anchor
composite-reward control (\texttt{use\_anchor=false}, $\lambda=0$;
otherwise identical hyperparameters; Table~\ref{tab:dstep-spread}).
Although this control does not use the anchor for selection, the
diagnostic can still be computed from its hidden-state trajectory.
\textsc{TADS} keeps both the displacement \emph{magnitude} (median
$0.068$, p95/median $1.4$) and the eigendirection \emph{sign} (zero
flips across $1{,}496$ measurements) uniformly tight. The no-anchor
control, sharing every detail except the anchor multiplier, produces
a $2.9\times$ larger median displacement, a single-layer max of
$1.998$ (a $5.8\times$ spike above \textsc{TADS}'s worst case), and
a sign-flip rate above the $5\%$ probabilistic-A3 threshold
(App.~F E3b $=$ FAIL). The point of contrast is not that the control
is uniformly worse on individual refreshes---by chance some
refreshes land inside \textsc{TADS}'s typical range---but that no
uniform guarantee is available: it is high-variance, occasionally
tight by accident and occasionally far outside.
% ----- TABLE: anchor displacement diagnostic -----
\begin{table}[t]
\centering
\caption{\textbf{Anchor-displacement diagnostic: \textsc{TADS} vs.\
matched no-anchor composite-reward control.} Same diagnostic, same
data, same hyperparameters except \texttt{use\_anchor} and $\lambda$.
The control does not use the anchor for selection, but its
hidden-state trajectory is measured post-hoc with the same
instrumentation. All values are post-warmup; sign-flip rate is per
layer-pair across 24 decoder layers and 68 refresh steps (max
possible: $24 \times 68 / 2 = 1{,}496$ pair-steps).}
\label{tab:dstep-spread}
\small
\setlength{\tabcolsep}{5pt}
\begin{tabular}{@{}lccccc@{}}
\toprule
 & \multicolumn{2}{c}{$d^{(t)}$, median over layers} & worst single & sign-flip & \\
\cmidrule(lr){2-3}
Cell & median & peak & layer-step & rate & E3b \\
\midrule
\textsc{TADS} (D)        & $0.068$ & $0.342$ & $0.342$ & $0 / 1{,}496$ & PASS \\
No-anchor control        & $0.199$ & $0.456$ & $1.998$ & ${>}\,5\%$    & FAIL \\
\bottomrule
\end{tabular}
\end{table}

\begin{figure}[t]
\centering
\resizebox{0.97\columnwidth}{!}{%
\begin{tikzpicture}[
  every node/.style={font=\scriptsize},
  >=Latex
]

% Plot area: x in [0, 9], y in [0, 6]. y is log10(value) mapped from
% [-3, +5] to [0, 6] via y_plot = (log10(value) + 3) * 0.75.
% Data: App. F cell B (probe 2000, cosine schedule, 78 refresh points
% every 25 optimizer steps from step 25 to step 1950, 24 decoder
% layers, Qwen2.5-0.5B + LoRA).
%   - red solid : measured cumulative drift to final anchor,
%                 median over layers of ||v_l^(t) - v_l^(T)||
%   - blue dash : Theorem 1 cumulative tail bound,
%                 (2 * C_sigma / gamma_min) * sum_{s >= t} eta_s

\path[use as bounding box] (-0.55,-0.95) rectangle (9.55,6.90);

% Light shading: region the bound covers (anywhere below the dashed
% line is theoretically safe).
\fill[blue!6] (0,0) -- (0.00,5.74) -- (0.48,5.73) -- (0.96,5.70)
              -- (1.44,5.66) -- (1.92,5.61) -- (2.40,5.57) -- (2.88,5.51)
              -- (3.36,5.45) -- (3.84,5.38) -- (4.32,5.30) -- (4.80,5.21)
              -- (5.28,5.11) -- (5.76,4.99) -- (6.24,4.86) -- (6.72,4.70)
              -- (7.20,4.50) -- (7.68,4.25) -- (8.16,3.91) -- (8.64,3.38)
              -- (9.00,2.60) -- (9.00,0) -- cycle;

% Y-axis grid lines at log10 = -3, -2, -1, 0, 1, 2, 3, 4, 5
\foreach \logv/\ylab in {-3/$10^{-3}$, -2/$10^{-2}$, -1/$10^{-1}$,
                          0/$10^{0}$,  1/$10^{1}$,   2/$10^{2}$,
                          3/$10^{3}$,  4/$10^{4}$,   5/$10^{5}$} {
  \pgfmathsetmacro{\ycoord}{(\logv + 3) * 0.75}
  \draw[gray!30, very thin] (0, \ycoord) -- (9, \ycoord);
  \node[anchor=east, font=\tiny] at (-0.05, \ycoord) {\ylab};
}

% X-axis tick labels (refresh step). 78 points span step 25..1950;
% x_plot = (step - 25) / 1875 * 9.
\foreach \step/\xc in {25/0.00, 500/2.28, 1000/4.68, 1500/7.08, 1950/9.00} {
  \draw[gray!50, very thin] (\xc, 0) -- (\xc, -0.08);
  \node[anchor=north, font=\tiny] at (\xc, -0.10) {\step};
}

% Theoretical cumulative tail bound (blue dashed). 20 sub-sampled points.
\draw[blue!70!black, very thick, dashed]
  plot[smooth] coordinates {
    (0.00,5.74) (0.48,5.73) (0.96,5.70) (1.44,5.66) (1.92,5.61)
    (2.40,5.57) (2.88,5.51) (3.36,5.45) (3.84,5.38) (4.32,5.30)
    (4.80,5.21) (5.28,5.11) (5.76,4.99) (6.24,4.86) (6.72,4.70)
    (7.20,4.50) (7.68,4.25) (8.16,3.91) (8.64,3.38) (9.00,2.60)
  };

% Measured cumulative drift to final anchor (red solid + dots).
\draw[red!75!black, thick]
  plot[smooth] coordinates {
    (0.00,2.09) (0.48,1.75) (0.96,1.73) (1.44,1.70) (1.92,1.68)
    (2.40,1.61) (2.88,1.59) (3.36,1.49) (3.84,1.41) (4.32,1.47)
    (4.80,1.36) (5.28,1.35) (5.76,1.33) (6.24,1.22) (6.72,1.07)
    (7.20,1.00) (7.68,0.77) (8.16,0.73) (8.64,0.64) (9.00,0.57)
  };
\foreach \p in {
  (0.00,2.09), (0.48,1.75), (0.96,1.73), (1.44,1.70), (1.92,1.68),
  (2.40,1.61), (2.88,1.59), (3.36,1.49), (3.84,1.41), (4.32,1.47),
  (4.80,1.36), (5.28,1.35), (5.76,1.33), (6.24,1.22), (6.72,1.07),
  (7.20,1.00), (7.68,0.77), (8.16,0.73), (8.64,0.64), (9.00,0.57)
} {
  \fill[red!75!black] \p circle (1.5pt);
}

% Axes drawn last
\draw[->, thick] (0,0) -- (9.25,0);
\draw[->, thick] (0,0) -- (0,6.40);

\node at (4.50,-0.55)
  {refresh step $t$ (cosine schedule, 78 points)};

\node[anchor=west, align=left, font=\scriptsize] at (-0.50,6.65)
  {\(\|v_l^{(t)} - v_l^{(T)}\|\) (median over 24 layers, log scale)};

% Legend (top-right, log-space y around 5.5--5.8)
\draw[blue!70!black, very thick, dashed] (5.20,6.30) -- (5.85,6.30);
\node[anchor=west, font=\tiny] at (5.95,6.30)
  {Thm.~1 tail bound: $\frac{2\hat{C}_\Sigma}{\gamma_{\min}}\sum_{s\ge t}\eta_s$};

\draw[red!75!black, thick] (5.20,5.85) -- (5.85,5.85);
\node[anchor=west, font=\tiny] at (5.95,5.85)
  {Measured $\|v^{(t)} - v^{(T)}\|$ (App.~F cell~B)};

\end{tikzpicture}%
}
\caption{\textbf{Empirical verification of
Theorem~\ref{thm:stability}: ε-stationarity under a cosine
schedule.} On Qwen2.5-0.5B + LoRA, App.~F cell~B (probe size
$2{,}000$, cosine learning-rate schedule, $78$ refresh points at
$25$-step intervals, $24$ decoder layers), we plot the cumulative
drift to the final anchor,
$\|v_l^{(t)} - v_l^{(T)}\|$ (red, median over layers), against the
Theorem~1 cumulative tail bound
$(2\hat{C}_\Sigma / \gamma_{\min}) \sum_{s\ge t}\eta_s$
(blue dashed; $\hat{C}_\Sigma = 5.78\times 10^{3}$,
$\gamma_{\min} = 0.050$). Both curves are shown on a log scale.
The bound is monotonically decreasing because the cosine schedule
makes the remaining-learning-rate tail $\sum_{s\ge t}\eta_s$ shrink
toward zero, so the bound itself goes to zero as $t \to T$; the
measured drift decays in lockstep, from $0.61$ at the first
refresh to $5.8\times 10^{-3}$ at the last. The measured drift
stays strictly below the bound at every refresh point and tracks
its shape rather than saturating, which is exactly the
$\varepsilon$-stationarity statement following
Eq.~\ref{eq:stability-bound}. The bound is loose but never vacuous
or violated; the no-anchor diagnostic of
Table~\ref{tab:dstep-spread} contrasts this with a matched control
in which the same diagnostic produces uncontrolled, unbounded
displacement.}
\label{fig:thm-main}
\end{figure}

\paragraph{Re-ranking overhead.}
\textsc{TADS} adds a structural signal without changing the
dominant compute terms of dynamic selection. The candidate-pool
forward pass and the fine-tuning update are shared with the
no-anchor composite-reward base. The only additional operations
are trajectory-anchor PCA and alignment scoring, both computed
from hidden states already produced during the forward pass.
Table~\ref{tab:cost-14b-main} shows these additions take about
$15$ seconds per epoch on Qwen2.5-14B, increasing the per-epoch
runtime by less than $0.5\%$.

\begin{table}[t]
\centering
\caption{\textbf{Per-epoch cost on Qwen2.5-14B + LoRA + Alpaca-GPT4}
($N=52{,}002$, $8\times$A100). The forward pass and fine-tuning
update are shared with the no-anchor composite-reward base. The
\textsc{TADS}-only cost comes from a small PCA and alignment dot
products over cached hidden states.}
\label{tab:cost-14b-main}
\small
\begin{tabular}{lcc}
\toprule
Stage & Composite-reward base & \textsc{TADS} \\
\midrule
Forward pass over pool          & $\sim$5 min  & $\sim$5 min  \\
Composite reward $\mathcal{R}_i^{(t)}$
                                & $\sim$3 s    & $\sim$3 s    \\
Trajectory anchor PCA           & ---          & $\sim$10 s   \\
Alignment score                 & ---          & $\sim$5 s    \\
Top-$B$ selection                & $\sim$1 s    & $\sim$1 s    \\
Fine-tuning on selected subset   & $\sim$30 min & $\sim$30 min \\
\midrule
Total per epoch                  & $\sim$35.1 min & $\sim$35.3 min \\
\textsc{TADS} overhead           & ---          & $<0.5\%$     \\
\bottomrule
\end{tabular}
\end{table}
